% !TEX root = ../notes_template.tex
\chapter{Knee}\label{chp:knee}

\section*{Objectives}

\begin{itemize}
  \item Analyze knee-relevant SFMA Top Tier patterns to identify dysfunction and regional interdependence.
  \item Perform MMT for knee extension and flexion, including assessment of primary and secondary contributors.
  \item Utilize hand-held dynamometry to evaluate quadriceps:hamstring strength ratios.
  \item Conduct the Nordic Hamstring Exercise as a screen for eccentric control.
  \item Perform and interpret the patellar tendon reflex to assess femoral nerve and L3–L4 spinal root integrity.
  \item Perform muscle length testing for the rectus femoris using the Modified Thomas Test and document R1/R2 values.
  \item Visualize the distal quadriceps tendon, patellar tendon, and assess suprapatellar effusion via RUSI.
  \item Palpate and image the popliteal pulse using ultrasound to assess vascular integrity.
\end{itemize}

\section*{Functional Integration: The Knee as a Force Transducer}

The knee functions as a dynamic hinge between the proximal power generator (hip) and the adaptable distal controller (foot/ankle complex). While it enables powerful sagittal plane motion, it is structurally vulnerable due to its relatively incongruent articular surfaces and heavy dependence on passive stabilizers (e.g., ACL, menisci). Clinically, knee symptoms often reflect issues above or below the joint, highlighting the importance of regional interdependence.

Efficient knee function depends on proximal control (hip extensors and rotators) and distal alignment (foot/ankle mechanics), as well as intact neuromuscular timing of the quadriceps and hamstrings.

Dysfunction at the knee often presents with symptoms locally but arises from regional interdependence above or below. For example, valgus collapse during gait or squat may reflect poor hip control or overpronation, while anterior knee pain may stem from quadriceps dominance, limited dorsiflexion, or impaired patellar tracking.

\vspace{0.25cm}
\subsection*{SFMA Top Tier Patterns with Knee Relevance}

\begin{itemize}
  \item \textbf{Multisegmental Flexion (MSF)} — Observe for incomplete knee extension during descent; restricted hamstring length or posterior capsule tightness may contribute to early spinal flexion.
  \item \textbf{Multisegmental Extension (MSE)} — Monitor for knee hyperextension or asymmetrical loading; anterior pelvic tilt or quadriceps dominance may elevate patellofemoral joint stress.
  \item \textbf{Multisegmental Rotation (MSR)} — Assess for tibiofemoral rotation asymmetries; impaired screw-home mechanics may result from femoral or foot/ankle dysfunction.
  \item \textbf{Single Leg Stance (SLS)} — Note frontal plane drift or valgus collapse at the knee; deficits in gluteal control, proprioception, or tibial stability may underlie poor stance control.
  \item \textbf{Arms Down Squat (ADS)} — Evaluate foot–knee–hip alignment, depth control, and valgus tendencies; limitations may suggest proximal motor deficits, distal mobility restrictions, or neuromechanical imbalance.
\end{itemize}

\noindent
This week’s lab draws from these patterns to evaluate knee motion, muscle function, and control using MMT, MLT, and RUSI. Use movement observations to guide focused assessment of neuromuscular and vascular systems.

\section*{Knee Movements: Muscles and Innervation}

\begin{table}[H]
\centering
\captionsetup{justification=raggedright, singlelinecheck=false}
\begin{tabular}{|p{4cm}|p{5cm}|p{4cm}|p{2.5cm}|}
\hline
\textbf{Movement} & \textbf{Muscle} & \textbf{Peripheral Nerve} & \textbf{Nerve Root(s)} \\
\hline
\multirow{2}{*}{Knee Extension}
  & Rectus Femoris & Femoral & L2--L4 \\
  & Vastus Group & Femoral & L2--L4 \\
\hline
\multirow{4}{*}{Knee Flexion}
  & Biceps Femoris (long head) & Sciatic (tibial division) & L5--S2 \\
  & Biceps Femoris (short head) & Sciatic (common fibular division) & L5--S2 \\
  & Semitendinosus & Sciatic (tibial division) & L5--S2 \\
  & Semimembranosus & Sciatic (tibial division) & L5--S2 \\
\hline
\end{tabular}
\caption{Primary Knee Flexors and Extensors with Neuromuscular Innervation.}
\end{table}

\subsection*{MMT and Performance Testing – Knee Flexion \& Extension}

\begin{labactivitybox}
\textbf{Activity: Manual Muscle Testing (MMT) – Knee Flexion and Extension}

\begin{itemize}
  \item Perform gravity-resisted and gravity-eliminated MMT for:
    \begin{itemize}
      \item \textbf{Knee Extension:} Primarily quadriceps group (femoral nerve, L2–L4)
      \item \textbf{Knee Flexion:} Primarily hamstrings (sciatic nerve, L5–S2)
    \end{itemize}
  \item Observe for compensations such as hip movement or foot rotation.
  \item Document muscle tested, innervation, and nerve root(s).
\end{itemize}
\end{labactivitybox}

\begin{labactivitybox}
\textbf{Activity: HHD Testing – Quadriceps and Hamstring Strength}

\textit{Quadriceps:}
\begin{itemize}
  \item Use HHD in seated knee extension position.
  \item Stabilize proximally and measure isometric force output.
\end{itemize}

\textit{Hamstrings:}
\begin{itemize}
  \item Use HHD in prone knee flexion position.
  \item Ensure neutral hip position and compare side-to-side.
\end{itemize}

\textbf{Record:}
\begin{itemize}
  \item Absolute and relative torque for each group.
  \item Calculate \textbf{H:Q ratio} to assess flexor/extensor balance.
\end{itemize}
\end{labactivitybox}

\begin{labactivitybox}
\textbf{Activity: Nordic Hamstring Exercise (NHE) – Eccentric Strength Evaluation}

\begin{itemize}
  \item Begin in tall kneeling with ankles stabilized by a partner.
  \item Lean forward slowly, controlling descent as long as possible.
  \item Emphasize slow, eccentric hamstring activation.
  \item Note the point of failure and control strategies.
  \item Optional: Use video or angle goniometer to document endpoint.
\end{itemize}

\textbf{Interpretation:}
\begin{itemize}
  \item Reflect on how torque relationships determine movement vs. control.
  \item Discuss implications for hamstring injury risk and athletic performance.
\end{itemize}

\textit{Reference:} Cuthbert et al. (2020)\footnote{The Nordic Hamstring Exercise has been shown to improve eccentric strength and increase fascicle length in the hamstrings, which may reduce injury risk. See Cuthbert et al., Sports Medicine, 50(1), 2020.}
\end{labactivitybox}

\subsection*{DTR: Patella}

\begin{labactivitybox}
\textbf{Activity: DTR – Patellar Tendon Reflex}

\begin{itemize}
  \item Position patient seated with knees flexed at ~90°.
  \item Tap patellar tendon using reflex hammer; observe quadriceps contraction.
  \item Grade on standard 0–4+ scale:
    \begin{itemize}
      \item 0 = Absent
      \item 1+ = Hypoactive
      \item 2+ = Normal
      \item 3+ = Brisk
      \item 4+ = Clonus
    \end{itemize}
  \item Correlate with femoral nerve integrity and L3–L4 spinal roots.
\end{itemize}
\end{labactivitybox}

\subsection*{MLT: Knee-Influencing Muscles}

\begin{labactivitybox}
\textbf{Activity: MLT -- Rectus Femoris (Ely's Test)}
\begin{itemize}
  \item With patient prone, passively flex the knee.
  \item Observe for anterior pelvic tilt or hip flexion as compensation.
  \item Document R1 and R2 and classify muscle length.
\end{itemize}
\end{labactivitybox}

\begin{labactivitybox}
\textbf{Activity: MLT -- Hamstrings (90-90 Test)}
\begin{itemize}
  \item Supine position with hips and knees at 90\degree; passively extend the knee.
  \item Record R1 and R2; observe for posterior pelvic tilt or lumbar flattening.
  \item Compare bilaterally and categorize extensibility.
\end{itemize}
\end{labactivitybox}


\section*{Rehabilitative Ultrasound Imaging of the Knee}

\subsection*{Musculotendonous Assessment and Patella Tracking}

\begin{labactivitybox}
\textbf{Activity: RUSI -- Patellar and Quadriceps Tendons}
\begin{itemize}
  \item Position patient supine or long sitting.
  \item Use SAX to identify tendons at the patellar apex and tibial tuberosity.
  \item Translate to track the distal quadriceps and the musculoskeltal junction
  \item Use LAX to follow fiber alignment.
  \item Have your partner perform quadriceps activation to see changes to the fiber alignment and patella movement
  \item Observe echogenicity, boundaries, and dynamic contraction.
\end{itemize}
\end{labactivitybox}

\subsection*{Knee Effusion Assessment}
\noindent
\textbf{Clinical Rationale: Knee Effusion Monitoring in Rehabilitation}

Effusion tracking is a powerful tool for load management in knee rehabilitation. Physical tests such as the bulge or sweep test have limited diagnostic precision: while they demonstrate a reasonably good \textit{positive likelihood ratio} (i.e., a positive finding increases the probability of a true effusion), they have a poor \textit{negative likelihood ratio}, meaning a negative finding does \textit{not} reliably rule out effusion. In contrast, ultrasound provides direct, quantifiable imaging of intra-articular fluid and enables consistent comparisons over time.

By routinely imaging the suprapatellar pouch and parapatellar recesses, clinicians can detect subtle changes in effusion volume that may indicate excessive tissue loading, delayed recovery from activity, or subclinical inflammation. Incorporating effusion RUSI into clinical decision-making supports better progression of exercise intensity, guides return-to-activity planning, and enhances long-term monitoring of joint status.
\vspace{0.5cm}

\begin{labactivitybox}
\textbf{Activity: RUSI Screening for Knee Effusion}

\textit{Effusion Evaluation – Suprapatellar Pouch and Parapatellar Recesses}

\begin{itemize}
  \item \textbf{Position:} Patient supine with the knee in slight flexion (20–30°), supported by a bolster.
  \item \textbf{Step 1 – Suprapatellar Pouch:}
    \begin{itemize}
      \item Use a \textbf{SAX} view transversely above the patella with the probe oriented horizontally.
      \item Identify the quadriceps tendon, femur, and fluid-filled suprapatellar space (if present).
      \item Look for anechoic (black) or hypoechoic (dark) fluid between the femur and soft tissues.
    \end{itemize}
  \item \textbf{Step 2 – Medial Parapatellar Recess:}
    \begin{itemize}
      \item Place the probe longitudinally just medial to the patella and distal to the joint line.
      \item Identify the joint capsule and look for anechoic distention or capsular bulging.
    \end{itemize}
  \item \textbf{Compare bilaterally} for side-to-side symmetry and document presence or absence of fluid.
  \item \textbf{Label and save images} using: \texttt{Rt/Lt SP Pouch} and \texttt{Rt/Lt MP Recess}.
\end{itemize}
\end{labactivitybox}


\section*{Pulse Palpation and Imaging}

\begin{labactivitybox}
\textbf{Activity: Popliteal Artery – Palpation and Ultrasound Imaging}

\begin{itemize}
  \item \textbf{Palpation:} Locate the popliteal artery with the patient supine and the knee slightly flexed. Use deep fingertip pressure at the midline of the popliteal fossa.
  \item \textbf{Ultrasound Imaging:}
    \begin{itemize}
      \item Use a low-frequency curvilinear or high-frequency linear probe (depending on patient size).
      \item Place the probe in \textbf{SAX} over the popliteal crease.
      \item Identify the artery as a pulsatile, non-compressible circular structure posterior to the femur and deep to the tibial nerve and vein.
      \item Apply light pressure to assess compressibility and confirm vessel identity.
      \item Optionally use color Doppler to visualize pulsatility.
    \end{itemize}
  \item \textbf{Clinical relevance:} Popliteal artery assessment can aid in peripheral vascular screening, especially in cases of posterior knee pain, swelling, or concern for vascular compromise.
\end{itemize}
\end{labactivitybox}

\section*{Clinical Integration}

\begin{itemize}
  \item How do knee impairments reflect dysfunction above (e.g., hip control) or below (e.g., foot alignment)? What SFMA Top Tier findings support this?
  \item How do functional strength measures (e.g., Nordic hamstring, HHD testing) enhance your interpretation of MMT findings?
  \item What role does the hamstring–quadriceps strength ratio play in injury risk and rehabilitation progression?
  \item How does ultrasound-based effusion assessment improve your ability to manage load and monitor tissue response over time?
  \item How did palpation and imaging of the popliteal artery reinforce anatomical reasoning and ultrasound identification skills?
\end{itemize}

\section*{Next Steps}

\begin{itemize}
  \item Practice MMT for knee flexion and extension, emphasizing stabilization and accurate resistance application.
  \item Review the application and interpretation of the Nordic hamstring test and HHD strength ratio analysis.
  \item Reinforce your ability to identify rectus femoris length limitations using the Modified Thomas Test.
  \item Practice scanning of the patellar tendon, distal quadriceps, and suprapatellar pouch with ultrasound.
  \item Preview Week 7’s ankle and foot content, focusing on medial arch mechanics, Achilles tendon imaging, and lateral ankle ligament scanning.
\end{itemize}

\section*{Checklist Summary}

\begin{itemize}
  \item [$\square$] Interpreted SFMA patterns (MSF, MSE, MSR, SLS, ADS) with attention to tibiofemoral mechanics and regional interdependence.
  \item [$\square$] Performed MMT for knee flexion and extension, identifying key muscle groups and neural innervation.
  \item [$\square$] Used hand-held dynamometry to compare quadriceps and hamstrings strength and analyze functional ratio.
  \item [$\square$] Performed the Nordic Hamstring Exercise as a screen for eccentric knee flexor control.
  \item [$\square$] Performed the patellar tendon reflex and interpreted findings in relation to L3–L4 and femoral nerve function.
  \item [$\square$] Conducted the Modified Thomas Test to evaluate rectus femoris length with R1 and R2 documentation.
  \item [$\square$] Acquired RUSI images of the distal quadriceps, patellar tendon, and suprapatellar pouch for effusion assessment.
  \item [$\square$] Palpated and imaged the popliteal artery, confirming pulsatility and compressibility with ultrasound.
\end{itemize}
